% %--------------------------
% %	CONFIGURATION
% %--------------------------
\documentclass[11pt,a4paper]{moderncv}
\moderncvstyle{classic}
\moderncvcolor{black}
\definecolor{firstnamecolor}{rgb}{0,0,0}

% Basic packages
\usepackage[utf8]{inputenc}
\usepackage[T1]{fontenc}
\usepackage[scale=0.7, top=2cm, bottom=2cm]{geometry}
\usepackage{microtype}

% Configure font settings
\renewcommand{\rmdefault}{ppl}
\renewcommand{\sfdefault}{phv}
\renewcommand*{\namefont}{\fontsize{26}{18}\mdseries}

% Configure microtype
\microtypesetup{expansion=false}

% Load sensitive information
\input{../../secrets/personal_info_dk.tex}

% %--------------------------
\begin{document}
\makecvtitle
Dear Members of the Selection Committee,

\vspace*{2em}
I am writing to apply for the PhD scholarship in Uncertainty Quantification for Large Language Models at DTU Compute. I have a Master's degree in Statistical Learning and Data Analytics from KTH, a background in Engineering Physics, and seven years of experience applying machine learning in research and production settings. My background is strongest in probabilistic modelling, model evaluation, and production ML systems, and I would like to develop that into research on how statistical uncertainty estimates can be calibrated, scaled, and communicated through LLM-generated text.

\hspace*{2em}
My strongest research experience comes from RaySearch Laboratories, where I worked on automated radiotherapy planning using deep learning, probabilistic modelling, and optimisation. There I built latent representations from segmented 3D CT scans with autoencoders and predicted dose-volume histograms and spatial dose distributions using sparse Gaussian processes. This taught me to work with uncertainty in a high-stakes setting and to treat evaluation as part of the modelling problem rather than an afterthought. It also gives me a practical basis for working with approximate Bayesian inference, calibration, and uncertainty-aware evaluation, which I see as central to making LLM confidence estimates meaningful rather than merely stylistic.

\hspace*{2em}
I have also worked in applied settings where model behaviour needed to be understood clearly. At Valcon, I built a production real-time neural network using NLP and high-dimensional labels, and I led explainable AI and Git trainings. At Signum Life Science, I currently work on a forecasting platform for pharmaceutical price and demand forecasting, and I use LLMs to embed free text for similarity search and matching workflows.

\hspace*{2em}
The topic also fits my academic background. At KTH, I specialised in Statistical Learning and Data Analytics, with a broader engineering foundation in mathematics and physics. Relevant coursework included Statistical Machine Learning, Probability Theory, and Reinforcement Learning. I also spent an exchange at ETH studying Applied Mathematics. That training is why I am comfortable moving from applied ML back towards mathematically grounded research. While my direct LLM work has mainly been on embeddings, retrieval, and evaluation rather than foundation-model training, my academic background in probabilistic modelling and production ML gives me a strong base for moving into uncertainty quantification for modern language models.

\hspace*{2em}
I am also attracted by the project's collaboration across DTU, ITU, and Edinburgh, since the problem sits naturally between probabilistic machine learning, NLP, and human-facing model evaluation. I would value the chance to study how uncertainty can be measured, calibrated, and communicated in LLMs in a way that is useful to researchers and end users alike.

\vspace*{2em}
I would welcome the opportunity to discuss my background and how it could contribute to the project. 

\vspace*{1em}
Regards, dIvar Cashin Eriksson.

\end{document}
